\chapter*{Preface to the archive}\label{preface}
\addcontentsline{toc}{chapter}{Preface to the archive}

This volume preserves the working record behind Bijou: philosophy and architecture, practical guides, design proposals, release and migration notes, retrospectives, audits and historical planning. It is the companion to \textit{The Bijou Book}, whose five Parts teach a reproducible job-monitor application.

\section*{What an archive can establish}

The 426 source records retain their paths and stable record labels. Their typeset text and illustrations are unchanged from the preceding illustrated edition. The original Markdown bytes remain in \path{source-corpus.json}. The source inventory connects each record with its original path and hash.

A status, target release or future-tense statement in a record describes that source at its original point in the work. This edition does not claim to have checked whether every proposal shipped, remains current or was superseded. Each record therefore carries an explicit historical notice. Follow the linked implementation, release and validation evidence before applying its claim to a current application.

\section*{How to use this volume}

Begin with the introduction to the relevant family. The design-record introduction explains the recurring document fields, including the difference between proposed tests and recorded evidence. The later releases, retrospectives and audits can help assess the consequences of an earlier choice. Historical METHOD, cycle and legend terminology explains the documents that used it; it is not a restored operating process.

The standalone archive and the complete edition identify records as R1 through R426. Page numbers differ between editions. The machine-readable page maps and original-path inventory identify the same record in each. The reference atlas documents its capture version separately from the published version taught by the companion application.
